\chapter{Virtual Memory}

\section{Physical vs virtual addresses}

\begin{itemize}
  \item \textbf{Physical memory}: if the machine has $P$ bytes of RAM, physical addresses range
    over $0 \ldots P-1$. A machine with 32-bit physical addresses can address at most
    $2^{32}$ bytes = 4\,GB.
  \item \textbf{Virtual memory}: the kernel gives each process its own \textbf{address space},
    a range of \textbf{virtual addresses} $0 \ldots 2^{V}-1$. All addresses that a process
    generates (instruction fetches, loads, stores) are virtual.
  \item The \textbf{MMU} (memory management unit) translates every virtual address to a
    physical address, using information set up by the kernel.
\end{itemize}

Why virtualize memory?
\begin{enumerate}
  \item \textbf{Isolation/protection}: a process can only reach physical memory that the
    kernel mapped into its address space.
  \item \textbf{Convenience}: every program can be linked to start at the same virtual address;
    the address space can be larger than physical memory.
  \item \textbf{Efficient use of memory}: only the parts actually used need physical memory
    (and pages can live on disk).
\end{enumerate}

\section{Dynamic relocation (base and bound)}

The simplest scheme: each address space occupies one contiguous region of physical memory.
The MMU has a \textbf{relocation register} $R$ (base) and a \textbf{limit register} $L$ (size).
\[
  \text{translate}(v) = \begin{cases} R + v & \text{if } v < L, \\ \text{exception} & \text{otherwise.}\end{cases}
\]
On a context switch the kernel reloads $R$ and $L$.

\begin{example}
$R = \mathtt{0x24000}$, $L = \mathtt{0x10000}$. Virtual \code{0x0F2A} $\to$ \code{0x24F2A}.
Virtual \code{0x10004} $\ge L$ $\to$ exception.
\end{example}

\begin{figure}[H]
\centering
\begin{tikzpicture}[x=1cm, y=1cm, font=\small, >=Stealth]
  % virtual
  \draw[thick, fill=blue!10] (0,0) rectangle (1.6,2);
  \node[font=\scriptsize, anchor=east] at (-0.1,0) {\texttt{0x0}};
  \node[font=\scriptsize, anchor=east] at (-0.1,2) {$L=$\,\texttt{0x10000}};
  \node[below] at (0.8,-0.2) {virtual};
  \draw[red!70!black, thick] (0,0.35) -- (1.6,0.35);
  \node[font=\scriptsize, red!70!black, anchor=east] at (-0.1,0.35) {\texttt{0x0F2A}};
  % physical
  \draw[thick] (5,-1) rectangle (7.2,4);
  \fill[gray!20] (5,-1) rectangle (7.2,0.2);
  \fill[gray!20] (5,2.6) rectangle (7.2,3.4);
  \draw[thick, fill=blue!10] (5,0.6) rectangle (7.2,2.6);
  \draw[thick] (5,-1) rectangle (7.2,4);
  \node[font=\scriptsize, anchor=west] at (7.3,0.6) {$R=$\,\texttt{0x24000}};
  \node[font=\scriptsize, anchor=west] at (7.3,2.6) {$R+L=$\,\texttt{0x34000}};
  \node[below] at (6.1,-1.2) {physical};
  \node[font=\scriptsize, gray] at (6.1,-0.4) {other process};
  \node[font=\scriptsize, gray] at (6.1,3.0) {other process};
  \draw[red!70!black, thick] (5,0.95) -- (7.2,0.95);
  \node[font=\scriptsize, red!70!black, anchor=west] at (7.3,0.95) {\texttt{0x24F2A}};
  \draw[->, dashed] (1.6,0) -- (5,0.6);
  \draw[->, dashed] (1.6,2) -- (5,2.6);
  \node[font=\scriptsize, align=center] at (3.3,1.7) {add $R$\\(if $v<L$)};
\end{tikzpicture}
\caption{Dynamic relocation: the whole address space is one contiguous
block of physical memory starting at $R$.}
\end{figure}

\textbf{Problem}: the whole address space (including the big unused gap between heap and stack)
must be physically contiguous. This causes \textbf{external fragmentation}: free physical memory
is split into holes too small to use.

\section{Segmentation}

Divide the address space into a few \textbf{segments} (code, data, heap, stack), each relocated
independently with its own base and limit. A virtual address is (segment number, offset).
With $K$ bits for the segment number:
\[
  v = \underbrace{s}_{\text{top }K\text{ bits}} \mid \underbrace{o}_{\text{offset}}, \qquad
  \text{translate}(v) = R_s + o \text{ if } o < L_s.
\]
Unused gaps between segments need no physical memory, but each segment is still contiguous,
so external fragmentation remains.

\begin{example}
32-bit virtual addresses, 2-bit segment number. Segment 1 has base \code{0x40000} and limit
\code{0x8000}. Virtual \code{0x40001F00} has segment $= 01_2 = 1$, offset \code{0x1F00} $<$ limit,
so physical $=$ \code{0x41F00}.
\end{example}

\begin{figure}[H]
\centering
\begin{tikzpicture}[x=1cm, y=0.5cm, font=\small, >=Stealth]
  % virtual space
  \draw[thick] (0,0) rectangle (1.8,8);
  \foreach \y/\h/\lab/\c in {0/1.2/code/blue!15, 1.2/1/data/green!15, 2.2/1.4/heap/orange!20, 6.6/1.4/stack/red!15}
    \draw[thick, fill=\c] (0,\y) rectangle node{\lab} (1.8,\y+\h);
  \node[gray, font=\scriptsize] at (0.9,4.4) {unused};
  \node[below] at (0.9,-0.2) {virtual};
  % physical
  \draw[thick] (5,0) rectangle (6.8,8);
  \foreach \y/\h/\lab/\c in {0.4/1.4/stack/red!15, 2.6/1.2/code/blue!15, 4.6/1.4/heap/orange!20, 6.6/1/data/green!15}
    \draw[thick, fill=\c] (5,\y) rectangle node{\lab} (6.8,\y+\h);
  \node[below] at (5.9,-0.2) {physical};
  \foreach \ay/\by in {0.6/3.2, 1.7/7.1, 2.9/5.3, 7.3/1.1} \draw[->, gray!70!black] (1.8,\ay) -- (5,\by);
\end{tikzpicture}
\caption{Segmentation: each segment is relocated independently, so the
unused gap between heap and stack needs no physical memory, but each
segment must still be contiguous.}
\end{figure}

\section{Paging}

\begin{definition}[Paging]
Physical memory is divided into fixed-size \textbf{frames}; each virtual address space is
divided into \textbf{pages} of the same size (typically 4\,KB). Any page can be placed in any
frame. A per-process \textbf{page table} maps page numbers to frame numbers.
\end{definition}

A virtual address with page size $2^{d}$ splits into a \textbf{page number} (high bits) and an
\textbf{offset} (low $d$ bits). The offset is unchanged by translation.

\begin{center}
\begin{tikzpicture}[b/.style={draw,minimum height=0.7cm,font=\small}]
  \node[b,minimum width=3cm,fill=blue!10] (vpn) {page number $p$};
  \node[b,minimum width=2cm,right=0pt of vpn,fill=green!10] (off) {offset $o$};
  \node[above=0.1cm of vpn.north west,anchor=south west,font=\small] {virtual address};
  \node[draw,below=1.4cm of vpn,minimum width=2.6cm,minimum height=2.6cm,align=center] (pt) {};
  \node[below=0pt of pt,font=\small] {page table};
  \foreach \i in {0,...,3} {
    \draw ($(pt.north west)+(0,-0.65*\i)$) -- ($(pt.north east)+(0,-0.65*\i)$);
  }
  \node[font=\small] at ($(pt.north)+(0.35,-1.62)$) {frame $f$ \, V};
  \draw[-{Stealth}] (vpn.south) -- node[left,font=\small]{index} ($(pt.north)+(-0.6,-1.62)$);
  \node[b,minimum width=3cm,fill=orange!15,right=2.5cm of pt.east,yshift=0.2cm] (pfn) {frame number $f$};
  \node[b,minimum width=2cm,right=0pt of pfn,fill=green!10] (off2) {offset $o$};
  \node[below=0.1cm of pfn.south west,anchor=north west,font=\small] {physical address};
  \draw[-{Stealth}] ($(pt.north east)+(0,-1.62)$) -- (pfn.west);
  \draw[-{Stealth}] (off.south) to[out=-60,in=90] (off2.north);
  \node[draw,dashed,font=\small,left=0.6cm of pt] (ptbr) {PTBR};
  \draw[-{Stealth}] (ptbr) -- (pt.west);
\end{tikzpicture}
\end{center}

Each \textbf{page table entry (PTE)} contains the frame number plus flag bits:
\begin{description}
  \item[Valid] the page is mapped; access to an invalid page raises an exception.
  \item[Write / read-only] protection (e.g.\ code pages read-only); also execute bits.
  \item[Referenced (use)] set by hardware on access; used by page replacement.
  \item[Dirty (modified)] set on write; a dirty page must be written back before eviction.
  \item[Present] the page is in memory (vs.\ on disk) --- see page faults.
\end{description}

\begin{example}[Translation]
16-bit virtual addresses, 4\,KB ($2^{12}$) pages: 4-bit page number, 12-bit offset.
Page table: page $0 \to$ frame \code{0x0F}, page $1 \to$ frame \code{0x26}, page $2$ invalid,
page $3 \to$ frame \code{0x01}.
\begin{itemize}
  \item \code{0x102C}: page 1, offset \code{0x02C} $\to$ \code{0x2602C}.
  \item \code{0x3FFF}: page 3, offset \code{0xFFF} $\to$ \code{0x01FFF}.
  \item \code{0x2004}: page 2 invalid $\to$ exception.
\end{itemize}
\end{example}

\begin{important}[Paging pros and cons]
No external fragmentation (all frames are interchangeable). Some \textbf{internal
fragmentation}: on average half a page is wasted at the end of each segment. Main cost: the page
table itself. With 32-bit addresses and 4\,KB pages there are $2^{20}$ pages; with 4-byte PTEs a
\emph{linear} page table is 4\,MB \emph{per process}, even if the process uses little memory.
With 48-bit or 64-bit virtual addresses, linear tables are impossible.
\end{important}

\section{Multi-level page tables}

Make the page table itself paged: split the page number into several pieces. The first piece
indexes a top-level \textbf{directory}; each directory entry points to a second-level page table
(or is invalid, in which case that whole region needs no table at all).

\begin{center}
\begin{tikzpicture}[b/.style={draw,minimum height=0.7cm,font=\small},t/.style={draw,minimum width=1.5cm,minimum height=2cm}]
  \node[b,minimum width=2cm,fill=blue!10] (p1) {$p_1$};
  \node[b,minimum width=2cm,right=0pt of p1,fill=blue!20] (p2) {$p_2$};
  \node[b,minimum width=2cm,right=0pt of p2,fill=green!10] (o) {$o$};
  \node[t,below left=1.2cm and 0.0cm of p1] (dir) {};
  \node[below=0pt of dir,font=\small] {directory};
  \node[t,right=1.6cm of dir] (pt2) {};
  \node[below=0pt of pt2,font=\small] {2nd-level table};
  \node[b,minimum width=2.2cm,right=1.6cm of pt2,fill=orange!15] (f) {frame $f$};
  \node[b,minimum width=1.4cm,right=0pt of f,fill=green!10] (o2) {$o$};
  \draw ($(dir.west)+(0,0.2)$) -- ($(dir.east)+(0,0.2)$);
  \draw ($(dir.west)+(0,-0.15)$) -- ($(dir.east)+(0,-0.15)$);
  \draw ($(pt2.west)+(0,-0.4)$) -- ($(pt2.east)+(0,-0.4)$);
  \draw ($(pt2.west)+(0,-0.75)$) -- ($(pt2.east)+(0,-0.75)$);
  \draw[-{Stealth}] (p1.south) -- ($(dir.north)+(0,-0.8)$);
  \draw[-{Stealth}] ($(dir.east)+(0,0.03)$) -- ($(pt2.north west)+(0,-0.1)$);
  \draw[-{Stealth}] (p2.south) -- ($(pt2.north)+(0,-1.4)$);
  \draw[-{Stealth}] ($(pt2.east)+(0,-0.58)$) -- (f.west);
  \draw[-{Stealth}] (o.south) to[out=-70,in=110] (o2.north);
\end{tikzpicture}
\end{center}

\begin{example}[Two-level table for 32-bit addresses]
4\,KB pages $\Rightarrow$ 12-bit offset, 20-bit page number split $10 + 10$. Each table has
$2^{10}$ 4-byte entries $=$ exactly one 4\,KB page. A small process with one code page, one data
page and one stack page near the top needs the directory plus 2 second-level tables
($\approx$12\,KB), instead of 4\,MB.

General rule: choose the bits per level so each table fits in one page. With $2^{d}$-byte pages
and $2^{e}$-byte PTEs, each level resolves $d-e$ bits, so the number of levels is
$\ceil{(V-d)/(d-e)}$. E.g.\ x86-64: $V=48$, $d=12$, $e=3$ $\Rightarrow$ $\ceil{36/9} = 4$ levels.
\end{example}

Cost: a $k$-level table needs $k$ memory accesses to translate one address. This is why a TLB is
essential.

\section{Translation lookaside buffer (TLB)}

\begin{definition}[TLB]
A \textbf{TLB} is a small, fast, fully-associative cache in the MMU holding recent
(page number $\to$ frame number, flags) translations. On a TLB \textbf{hit}, translation takes no
extra memory access. On a \textbf{miss}, the translation must be found in the page table and
loaded into the TLB.
\end{definition}

\begin{itemize}
  \item \textbf{Hardware-managed TLB} (x86): on a miss the MMU walks the page table itself; the
    format of the page table is fixed by the hardware.
  \item \textbf{Software-managed TLB} (MIPS): on a miss the MMU raises an exception; the
    kernel's handler finds the translation (any data structure it likes) and writes it into the
    TLB with \code{tlbwr} (random slot) or \code{tlbwi} (given index).
\end{itemize}

Effective access time with hit rate $h$, memory access time $m$, TLB lookup negligible, and a
$k$-level table walk on a miss:
\[ \text{EAT} = h\cdot m + (1-h)(k+1)m. \]
E.g.\ $h = 0.99$, $k=2$: EAT $= 0.99m + 0.03m = 1.02m$. Programs have \emph{locality}, so
hit rates are typically very high.

On a context switch, TLB entries of the old process are wrong for the new one: either
\textbf{flush} (invalidate) the TLB, or tag each entry with an \textbf{address space ID}
(ASID). The MIPS R3000 TLB has 64 entries, each with a VPN in \code{EntryHi} (with a 6-bit ASID
that OS/161 ignores) and PFN plus \emph{valid}, \emph{dirty} (= writable), \emph{global},
\emph{nocache} bits in \code{EntryLo}. OS/161's \code{as\_activate} simply invalidates all TLB
entries.

\section{OS/161 dumbvm}

OS/161 starts with \code{dumbvm}, a deliberately simple VM system:
\begin{itemize}
  \item Each address space has exactly two segments (text and data) and a fixed-size stack
    (12 pages), each \emph{contiguous in physical memory}: essentially segmentation with base
    and size stored in \code{struct addrspace} (\code{as\_vbase1}, \code{as\_pbase1},
    \code{as\_npages1}, \ldots).
  \item Physical memory is grabbed with \code{ram\_stealmem} and never freed.
  \item On a TLB miss, \code{vm\_fault} computes the physical address by relocation and writes a
    TLB entry.
\end{itemize}
\begin{lstlisting}[style=c]
int vm_fault(int faulttype, vaddr_t faultaddress) {
    struct addrspace *as = proc_getas(curproc);
    vaddr_t vbase1 = as->as_vbase1, vtop1 = vbase1 + as->as_npages1 * PAGE_SIZE;
    vaddr_t vbase2 = as->as_vbase2, vtop2 = vbase2 + as->as_npages2 * PAGE_SIZE;
    vaddr_t stackbase = USERSTACK - DUMBVM_STACKPAGES * PAGE_SIZE;
    paddr_t paddr;

    faultaddress &= PAGE_FRAME;               /* page-align */
    if (faulttype == VM_FAULT_READONLY) return EFAULT;  /* A3: kill process */
    if (faultaddress >= vbase1 && faultaddress < vtop1)
        paddr = faultaddress - vbase1 + as->as_pbase1;
    else if (faultaddress >= vbase2 && faultaddress < vtop2)
        paddr = faultaddress - vbase2 + as->as_pbase2;
    else if (faultaddress >= stackbase && faultaddress < USERSTACK)
        paddr = faultaddress - stackbase + as->as_stackpbase;
    else
        return EFAULT;                        /* segmentation fault */

    int spl = splhigh();
    uint32_t ehi = faultaddress;
    uint32_t elo = paddr | TLBLO_DIRTY | TLBLO_VALID;
    tlb_random(ehi, elo);                     /* A3: replace a random entry when full */
    splx(spl);
    return 0;
}
\end{lstlisting}
Assignment 3 improves this: make the text segment read-only after loading (clear
\code{TLBLO\_DIRTY}), handle a full TLB, and replace \code{ram\_stealmem} with a \textbf{coremap}
that tracks which physical frames are free so memory can be reclaimed.

\section{Where address spaces come from}
The kernel builds an address space from an ELF executable: the ELF header lists
\emph{segments} (virtual address, size in file, size in memory, permissions). For each, the
kernel defines a region (\code{as\_define\_region}), allocates memory (\code{as\_prepare\_load}),
copies the file image in, zero-fills the rest (\code{.bss}), and finally sets up the stack
(\code{as\_define\_stack}). In OS/161 user stacks grow down from \code{USERSTACK =
0x80000000}.

\section{Paging to secondary storage}

Physical memory can be smaller than the sum of all address spaces. Keep only some pages in RAM;
the rest live in a \textbf{swap} area on disk. The PTE \emph{present} bit says which.

\begin{definition}[Page fault]
A \textbf{page fault} is the exception raised when a process accesses a valid page that is not
present in memory. The kernel:
\begin{enumerate}
  \item finds a free frame, or evicts a victim page (writing it to disk if dirty, and marking
    its PTE not present);
  \item reads the needed page from disk into the frame (the faulting thread blocks; another
    thread runs meanwhile);
  \item updates the PTE (present, frame number) and returns, re-executing the faulting
    instruction.
\end{enumerate}
\end{definition}

Demand paging: pages are loaded only when first touched. Prefetching (read nearby pages
together) can reduce fault counts.

\begin{important}[Cost]
A disk access is $\sim$5--10\,ms vs.\ $\sim$100\,ns for memory: a factor of $10^{5}$. With fault
rate $p$, EAT $= (1-p)\cdot 100\,\text{ns} + p \cdot 10\,\text{ms}$. To keep the slowdown under
10\%, $p$ must be below about $10^{-6}$. Replacement policy matters.
\end{important}

\section{Page replacement policies}

Measured with a \textbf{reference string} (sequence of page numbers) and a number of frames.

\subsection{Optimal (MIN / Belady)}
Evict the page whose next use is furthest in the future. Provably minimal faults but
requires knowledge of the future; used as a benchmark.

\subsection{FIFO}
Evict the page that has been in memory longest. Simple but ignores usage.

\begin{example}[3 frames, reference string \code{a b c d a b e a b c d e}]
\begin{center}
\begin{tabular}{@{}l*{12}{c}@{}}
\toprule
Ref & a & b & c & d & a & b & e & a & b & c & d & e \\
\midrule
FIFO F1 & a & a & a & d & d & d & e & e & e & e & e & e \\
FIFO F2 &   & b & b & b & a & a & a & a & a & c & c & c \\
FIFO F3 &   &   & c & c & c & b & b & b & b & b & d & d \\
Fault & x & x & x & x & x & x & x &   &   & x & x &   \\
\bottomrule
\end{tabular}
\end{center}
FIFO: 9 faults. OPT on the same string: 7 faults (a, b, c, d, e, c, d). With 4 frames FIFO
gives 10 faults --- more frames, more faults! This is \textbf{Belady's anomaly}.
\end{example}

\subsection{LRU (least recently used)}
Evict the page whose most recent use is furthest in the past --- past behaviour predicts the
future (locality). LRU is a \emph{stack algorithm} and does not suffer Belady's anomaly. On the
string above with 3 frames LRU gives 10 faults (this string is adversarial for it).
Exact LRU requires updating a timestamp or list on \emph{every} memory reference, which is too
expensive; hardware only provides the use bit.

\subsection{Clock (second chance)}
Approximates LRU with the \emph{use bit}. Frames are arranged in a circle with a ``hand''.
\begin{lstlisting}[style=c]
int clock_choose_victim(void) {
    for (;;) {
        struct frame *f = &frames[hand];
        hand = (hand + 1) % nframes;
        if (f->use) f->use = 0;     /* give a second chance */
        else return f - frames;     /* not used since last sweep: evict */
    }
}
\end{lstlisting}
Hardware sets \code{use = 1} on every access. On a MIPS software-managed TLB, the kernel can
simulate use bits by invalidating TLB entries and setting the bit when the next TLB miss for
that page occurs. A common refinement also prefers clean pages (use bit, dirty bit) over dirty
ones.

\begin{figure}[H]
\centering
\begin{tikzpicture}[font=\small, >=Stealth]
  \foreach \i/\pg/\u/\c in {0/A/1/blue!12, 1/B/0/white, 2/C/1/blue!12, 3/D/1/blue!12, 4/E/0/white, 5/F/1/blue!12, 6/G/0/white, 7/H/1/blue!12} {
    \pgfmathsetmacro{\a}{90-45*\i}
    \node[draw, thick, rounded corners, minimum width=1.1cm, minimum height=0.7cm, fill=\c] (f\i) at (\a:2.1) {\pg\ \textcolor{gray}{\scriptsize u=\u}};
  }
  \draw[->, very thick, red!70!black] (0,0) -- (90:1.5);
  \fill (0,0) circle (2pt);
  \node[font=\scriptsize, red!70!black] at (0.1,-0.35) {hand};
  \draw[->, thick, gray] (110:2.95) arc (110:70:2.95);
\end{tikzpicture}
\caption{The clock algorithm. The hand is at $A$ ($u=1$): it clears $A$'s
use bit and moves on; $B$ has $u=0$, so $B$ is evicted.}
\end{figure}

\subsection{Working sets and thrashing}
\begin{definition}
The \textbf{working set} $W(t,\Delta)$ of a process is the set of pages it referenced in the last
$\Delta$ time units (or references). \textbf{Thrashing} occurs when the sum of working sets
exceeds physical memory: processes fault continually, the CPU is mostly idle waiting for the
disk, and throughput collapses.
\end{definition}
Remedies: give each process enough frames for its working set (\emph{page-fault frequency}
control: if a process faults too often give it more frames, if rarely take some away), and if
memory is overcommitted, suspend (swap out) whole processes to reduce the
\emph{multiprogramming level}.

\begin{figure}[H]
\centering
\begin{tikzpicture}[x=1cm, y=1cm, font=\small, >=Stealth]
  \draw[->] (0,0) -- (7.4,0) node[below left, font=\scriptsize, align=right]{degree of multiprogramming};
  \draw[->] (0,0) -- (0,3.3) node[above, font=\scriptsize]{CPU utilization};
  \draw[very thick, blue!60!black] plot[smooth, tension=0.6] coordinates {(0,0) (1,1.4) (2.2,2.3) (3.6,2.75) (4.5,2.6) (5.2,1.6) (6,0.6) (7,0.3)};
  \draw[densely dashed, gray] (4.0,0) -- (4.0,2.75);
  \node[font=\scriptsize, red!70!black, align=center] at (6.0,1.6) {thrashing};
\end{tikzpicture}
\caption{Adding processes raises utilization until the working sets no
longer fit in memory; beyond that point the system thrashes.}
\end{figure}

\begin{aside}
Locality: \emph{temporal} (recently used data is used again soon) and \emph{spatial} (nearby
addresses are used soon). Both are what make TLBs, caches and LRU-like replacement effective.
\end{aside}

\section*{Summary}
\begin{itemize}
  \item The MMU translates virtual to physical addresses on every access; the kernel sets up
    the mapping and handles exceptions.
  \item Base/bound and segmentation are simple but suffer external fragmentation; paging avoids
    it at the cost of page tables (internal fragmentation is small).
  \item Multi-level page tables save space for sparse address spaces; TLBs make translation
    fast. MIPS uses a software-managed 64-entry TLB.
  \item \code{dumbvm} = contiguous segments + TLB refill by relocation.
  \item Page faults bring pages in from disk; replacement policies: OPT (ideal), FIFO
    (Belady's anomaly), LRU (good, expensive), clock (practical approximation).
  \item Thrashing occurs when working sets do not fit in memory.
\end{itemize}
